\documentclass[11pt]{article}

% ── Page geometry ──────────────────────────────────────────────────────────────
\usepackage[a4paper, top=25mm, bottom=28mm, left=20mm, right=20mm]{geometry}

% ── Fonts ──────────────────────────────────────────────────────────────────────
\usepackage[T1]{fontenc}
\usepackage{mathpazo}          % Palatino for body text and math
\usepackage{microtype}

% ── Math / science ─────────────────────────────────────────────────────────────
\usepackage{amsmath}
\usepackage{amssymb}
\usepackage{siunitx}

% ── Layout ─────────────────────────────────────────────────────────────────────
\usepackage{parskip}
\usepackage{enumitem}
\usepackage{setspace}

% ── Color (loaded before hyperref and colortbl) ────────────────────────────────
\usepackage[dvipsnames]{xcolor}
\definecolor{accent}{HTML}{1B4F8A}       % deep navy blue
\definecolor{accentlight}{HTML}{EBF2FA}  % pale blue for table header rows
\definecolor{codebg}{HTML}{F4F6F9}       % near-white blue-tinted code background
\definecolor{codenums}{HTML}{A0B4C8}     % muted blue-gray for line numbers

% ── Figures / tables ───────────────────────────────────────────────────────────
\usepackage{graphicx}
\usepackage{float}
\usepackage{booktabs}
\usepackage{array}
\usepackage{colortbl}
\usepackage[labelfont=bf, font=small]{caption}

% ── Code – tcolorbox with listings ────────────────────────────────────────────
\usepackage{listings}
\usepackage[breakable,skins,listings]{tcolorbox}

% C++ code style (used inside cppcode environment)
\lstdefinestyle{cppinner}{
  language=C++,
  basicstyle=\ttfamily\footnotesize,
  keywordstyle=\color{accent}\bfseries,
  commentstyle=\color{gray!70}\itshape,
  stringstyle=\color{BurntOrange},
  numbers=left,
  numberstyle=\tiny\color{codenums},
  numbersep=10pt,
  tabsize=4,
  breaklines=true,
  showstringspaces=false,
  keepspaces=true
}

% Plain-text format style (used inside fmtcode environment)
\lstdefinestyle{fmtinner}{
  language={},
  basicstyle=\ttfamily\small,
  numbers=none,
  tabsize=4,
  breaklines=true,
  showstringspaces=false
}

% Styled C++ code block: flat bg, accent left rule, no frame border
\newtcblisting{cppcode}[1][]{
  enhanced, breakable,
  arc=0pt, outer arc=0pt,
  colback=codebg,
  colframe=codebg,
  boxrule=0pt,
  borderline west={3pt}{0pt}{accent},
  top=5pt, bottom=5pt, left=4pt, right=4pt,
  listing only,
  listing options={style=cppinner},
  #1
}

% Styled plain-text format block: neutral bg, gray left rule
\newtcblisting{fmtcode}[1][]{
  enhanced, breakable,
  arc=0pt, outer arc=0pt,
  colback=gray!5,
  colframe=gray!5,
  boxrule=0pt,
  borderline west={3pt}{0pt}{gray!45},
  top=5pt, bottom=5pt, left=4pt, right=4pt,
  listing only,
  listing options={style=fmtinner},
  #1
}

% ── Headers / footers ──────────────────────────────────────────────────────────
\usepackage{fancyhdr}
\pagestyle{fancy}
\fancyhf{}
\fancyhead[L]{\small\color{gray!65} Sand Dune \textperiodcentered\ BOS System Documentation}
\fancyhead[R]{\small\color{gray!65} \leftmark}
\fancyfoot[C]{\small\color{gray!65} \thepage}
\renewcommand{\headrulewidth}{0.3pt}

% ── Section formatting ─────────────────────────────────────────────────────────
\usepackage{titlesec}
\titleformat{\section}
  {\large\bfseries\color{accent}}
  {\thesection}{1em}{}
  [{\vspace{2pt}\color{accent!30}\hrule height 0.6pt\vspace{4pt}}]
\titleformat{\subsection}
  {\normalsize\bfseries}
  {\thesubsection}{1em}{}
\titleformat{\subsubsection}
  {\normalsize\itshape}
  {\thesubsubsection}{1em}{}
\titlespacing{\section}{0pt}{2.5ex}{1.2ex}
\titlespacing{\subsection}{0pt}{2ex}{0.8ex}

% ── Links ──────────────────────────────────────────────────────────────────────
\usepackage[colorlinks=true, linkcolor=accent, citecolor=accent!80!black,
            urlcolor=accent!80!black]{hyperref}
\usepackage{xurl}
\urlstyle{tt}

% ── Custom commands ────────────────────────────────────────────────────────────
% Inline code: accent-colored monospace, with line-break permission after
% common delimiters so long identifiers don't overflow the right margin.
\newcommand{\codefmt}[1]{%
  \begingroup
  \ttfamily
  \hyphenchar\font=`\-
  \emergencystretch=3em
  \hspace{0pt}#1\hspace{0pt}%
  \endgroup}
\newcommand{\code}[1]{{\color{accent!85!black}\codefmt{#1}}}
\newcommand{\file}[1]{{\color{accent!70!black}\codefmt{#1}}}
\newcommand{\cls}[1]{{\bfseries\color{accent}\codefmt{#1}}}
\newcommand{\param}[1]{{\color{accent!60!black}\codefmt{#1}}}

% ── Figure borders ─────────────────────────────────────────────────────────────
% Use \figimg instead of \includegraphics to get a tight black border.
\newcommand{\figimg}[2][]{%
  \tcbox[enhanced, arc=4pt,
         boxrule=4pt, colframe=black, colback=white,
         boxsep=0pt, left=0pt, right=0pt, top=0pt, bottom=0pt]{%
    \includegraphics[#1]{#2}}%
}

% ═══════════════════════════════════════════════════════════════════════════════
\begin{document}

\begin{titlepage}
  \centering
  \vspace*{4cm}
  {\Huge\bfseries Sand Dune\par}
  \vspace{0.5em}
  {\large Design, Modeling, and Validation of a\\
   Background Oriented Schlieren System \\
   for High Resolution Refractive Index Mapping \\
   of Transparent Media\par}
  \vspace{2em}
  \hrule
  \vspace{1.5em}
  {\large Xander Linzel\\
   \href{mailto:xander.linzel@irradiant.tech}{\texttt{xander.linzel@irradiant.tech}}\par}
  \vspace{1em}
  {\normalsize \today\par}
  \vfill
  \begin{abstract}
    \setlength{\parindent}{0pt}
    Sand Dune is a custom Background Oriented Schlieren (BOS) system for
    quantitative measurement of weak refractive disturbances in transparent
    media. The project combines a dedicated optical rig with a desktop
    application that handles image loading, windowed cross-correlation,
    vector validation, surface reconstruction, visualization, and export.
    This document describes the optical geometry and physical model, the
    numerical processing pipeline in detail, how to use the application,
    and the testing strategy used to validate the software during development.
  \end{abstract}
\end{titlepage}

\tableofcontents
\clearpage

% ═══════════════════════════════════════════════════════════════════════════════
\section{Introduction}
\label{sec:intro}

Background Oriented Schlieren converts transverse refractive index gradients into an
image-registration problem. A structured background pattern is imaged through
a transparent specimen, and local refractive-index gradients cause apparent
displacements of the pattern in the recorded image. Comparing a reference
image (no specimen) with a disturbed image (specimen present) and measuring
those pixel shifts gives the gradient field of the refractive-index distribution.
Integrating those gradients yields a relative surface map of either the
refractive-index variation or the thickness variation of the sample,
depending on which quantity is assumed known.

BOS is attractive for weak optical disturbances, the kind where
interferometry would be impractical and shadowgraphy lacks quantitative
resolution. However, the accuracy of the final result depends on the entire
measurement chain: background pattern design, optical geometry, image quality,
sub-pixel estimation, validation, and reconstruction. Sand Dune was built to
address that full problem as a single coherent system rather than as a
stand-alone post-processing utility.

This document is organized as follows.
Section~\ref{sec:system} describes the physical setup.
Section~\ref{sec:model} presents the BOS optical model and the scaling
equations implemented in the code.
Section~\ref{sec:pipeline} describes each stage of the processing pipeline in
detail, including the cross-correlator, CMM sub-pixel solver, optional PID
deformation correction, validation, and surface reconstruction.
Section~\ref{sec:software} documents the codebase organization and key classes.
Section~\ref{sec:usage} explains how to use the application.
Section~\ref{sec:validation} summarizes the validation strategy and presents
the development-verification precision figures used to characterize the
current implementation.
Section~\ref{sec:precision} derives a minimum resolvable gradient for
the current optical configuration and compares it to the expected
measurement range.
Section~\ref{sec:calibration} discusses practical calibration.
Section~\ref{sec:caveats} lists known caveats in the current implementation.

% ═══════════════════════════════════════════════════════════════════════════════
\section{System Design}
\label{sec:system}

\subsection{Measurement Objective}

Sand Dune targets weak to moderate refractive disturbances in transparent
samples, a regime where displacements are often sub-pixel, smooth, and
low in magnitude. This makes the measurement sensitive to small errors at
every stage: optical alignment, background contrast, sub-pixel fitting, and
reconstruction. The physical system was therefore designed around three
practical constraints:

\begin{itemize}[noitemsep]
  \item stable optical alignment between camera, specimen, and background,
  \item repeatable sample positioning,
  \item a background pattern with sufficiently dense local texture for
        reliable sub-pixel cross-correlation.
\end{itemize}

\subsection{Optical Layout}

In the assembled rig, a camera with a fixed focal-length lens images a
patterned background target. The specimen or flow field is placed between
the background and the lens. The four geometric parameters that define this
layout are $Z_d$, $Z_a$, $f$, and $P_{px}$. These quantities are introduced
formally in Section~\ref{sec:params}; here it is sufficient to note that
they correspond to measurable hardware properties and inter-component
distances set during assembly and calibration.

The system is closely aligned with the BOS geometry described by
Raffel~\cite{raffel2015bos} and with later work on weak-flow extraction
using dense background texture and careful geometric design~\cite{liu2022uavbos}.

\subsection{Physical Hardware}

The hardware consists of a camera, a fixed focal-length imaging lens, a
patterned background target, a rigid support structure, and a specimen region
between the background and the lens. The region of interest may occupy only
part of the image frame, so the software supports a circular mask that
restricts the reconstruction domain to the sample area and excludes the
holder and other irrelevant structure. A custom holder was designed so that
the specimen position is repeatable and the masked reconstruction domain
matches the useful optical aperture of the setup.

\begin{figure}[H]
  \centering
  \figimg[width=0.5\linewidth]{hardware_cad.jpg}
  \caption{CAD design of Sand Dune optical rig showing camera, lens, specimen region,
           and patterned background target.}
  \label{fig:hardware}
\end{figure}

\begin{figure}[H]
  \centering
  \figimg[width=0.5\linewidth]{hardware_rig.jpg}
  \caption{Assembled Sand Dune optical rig showing camera, lens, specimen region,
           and patterned background target.}
  \label{fig:background}
\end{figure}

\subsubsection{Illumination}

The background is illuminated by an LED backlight panel operated at full
brightness. An incoherent light source is essential: coherent sources
(lasers) produce speckle in the image plane, which adds structured noise
to the cross-correlation and degrades sub-pixel accuracy. LEDs are suitable
because they are spatially incoherent.

A blue optical filter is placed in the illumination path. White-light
sources produce chromatic aberration,
which manifests as colour fringing around high-contrast
dot edges and introduces a small but systematic bias into the correlation
displacement estimate. The blue filter narrows the effective spectral
bandwidth, reducing chromatic aberration without requiring an apochromatic
lens.

\subsubsection{Background Patterns}

Six printed background patterns were fabricated, covering two dot densities
(40\,\% and 60\,\% fill) and three physical dot diameters (0.10\,mm,
0.50\,mm, and 0.75\,mm). The pattern is chosen to match the imaging
magnification so that each dot projects to approximately 4--5 pixels on
the sensor. This size is the practical optimum for FFT-based cross-correlation:

\begin{itemize}[noitemsep]
  \item Dots smaller than $\sim3$\,px are under-sampled; the correlation
        peak is narrow and noisy, and the CMM bicubic fit has insufficient
        support to resolve sub-pixel offsets accurately.
  \item Dots larger than $\sim7$\,px reduce the number of independent
        features per interrogation window, lowering the signal-to-noise
        ratio of the correlation peak.
  \item At 4--5\,px the peak is well-resolved, the CMM objective has a
        smooth, broad basin, and the S2N ratio is typically maximised.
\end{itemize}

Dot focus is equally important. An out-of-focus pattern blurs the dots to
larger apparent diameters and reduces edge contrast. Both effects degrade
the correlation peak shape, which CMM relies on for sub-pixel mapping.
The background should be focused as sharply as the camera optics allow.

The 40\,\% density patterns provide more inter-dot spacing and are better
suited to large interrogation windows or coarser magnification. The 60\,\%
density patterns pack more texture per window and are preferred when small
window sizes are used or when the dot diameter is near the lower end of the
4--5\,px range. For the present optical configuration, the 0.10\,mm,
40\,\% density pattern provides the most reliable correlation quality. For
future configurations, the background should be re-selected after checking
the projected dot size and the resulting correlation quality with the
software.

\subsubsection{Imaging Lens and Upgrade Path}

The current setup uses a \textbf{25\,mm fixed focal-length lens}. This
provides a field of view and working distance compatible with the present
sample size and mechanical layout. A longer focal-length lens remains a
reasonable design option for future versions of the rig, but in BOS the
consequences are best understood in terms of spatial specificity and the
geometry required to preserve a usable field of view.

At fixed $Z_a$, increasing $f$ enlarges the background image on the sensor,
so each interrogation window corresponds to a smaller physical region in the
sample plane. In BOS terms this makes the measurement more spatially
specific: a localized refractive disturbance influences a narrower family of
rays, which reduces lateral cross-talk between neighbouring regions and can
make fine optical structure easier to isolate. The practical tradeoffs are:

\begin{itemize}
  \item \textbf{Improved locality of the BOS response.} Because each pixel
        subtends a narrower bundle of rays at the sample plane, refractive
        features that are close together are less likely to contaminate the
        same interrogation windows. This can improve the lateral specificity
        of both the gradient fields and the reconstructed surface.

  \item \textbf{Background re-matching becomes necessary.} At the same
        stand-off distance, the background dots project to more pixels, so
        the physical dot diameter usually has to be reduced to remain in the
        preferred 4--5\,px range. A longer lens therefore normally requires
        a different background pattern even if the rest of the rig is
        unchanged.

  \item \textbf{Field coverage becomes the primary geometric constraint.}
        If the same sample area must remain visible, the camera commonly has
        to be moved farther from the specimen and background. That change
        modifies $Z_a$, $Z_B$, and the sample-plane grid spacing, so the
        final BOS sensitivity is determined by the full reconfigured
        geometry rather than by focal length alone.

  \item \textbf{Sensitivity and specificity no longer move together
        automatically.} A longer lens can improve locality of the measured
        disturbance while still forcing a larger stand-off distance to keep
        the full sample in view. Depending on how the rig is re-scaled, that
        can either help or partially offset the increase in measurable
        pixel displacement.

  \item \textbf{Focus tolerance becomes tighter.} Longer lenses have
        shallower depth of field at a given aperture and working distance,
        so the background pattern and specimen surfaces must be positioned
        more carefully to avoid differential blur across the optical path.

  \item \textbf{Lens quality matters more as magnification increases.}
        Any substitution should be evaluated experimentally for contrast,
        aberration, and repeatability with the intended background pattern,
        because those effects directly influence the correlation peak shape
        used by the sub-pixel solver.
\end{itemize}

% ═══════════════════════════════════════════════════════════════════════════════
\section{BOS Measurement Model}
\label{sec:model}

\subsection{Optical Parameters}
\label{sec:params}

The software uses the following optical parameters, all configurable through
the \textbf{Parameters} panel:

\begin{center}
\begin{tabular}{lll}
  \toprule
  \rowcolor{accentlight} Symbol & Parameter name in UI & Description \\
  \midrule
  $P_{px}$ & Pixel pitch & Sensor pixel pitch (\si{\micro\metre/pixel}) \\
  $f$      & Focal length & Lens focal length (\si{mm}) \\
  $Z_d$    & $Z_d$ & Background-to-sample distance (\si{mm}) \\
  $Z_a$    & $Z_a$ & Sample-to-lens front principal plane distance (\si{mm}) \\
  $Z_B$    & N/A & $Z_B = Z_a + Z_d$, background-to-lens distance (\si{mm}) \\
  $t$      & Thickness & Sample thickness (\si{mm}) \\
  $n$      & Index & Sample refractive index (dimensionless) \\
  \bottomrule
\end{tabular}
\end{center}

The image distance for the background plane follows the thin-lens relation:
\begin{equation}
  z_i = \frac{f Z_B}{Z_B - f}.
  \label{eq:zi}
\end{equation}
The lateral magnification of the background plane onto the sensor is
\begin{equation}
  m_1 = \frac{z_i}{Z_B}.
  \label{eq:mag}
\end{equation}

\subsection{What BOS Actually Measures}
\label{sec:opl}

BOS fundamentally measures the gradient of the \emph{optical path length}
(OPL) through the sample, not the refractive index or the thickness
independently. The optical path length through a sample at lateral position
$(x, y)$ is
\begin{equation}
  \mathrm{OPL}(x,y) = \int_0^t n(x, y, z)\,dz,
\end{equation}
where the integral runs along the optical axis through the full sample depth.
For a sample that is thin relative to $Z_a$ and optically homogeneous in
depth but varying laterally, this reduces to $\mathrm{OPL}(x,y) \approx
n(x,y)\,t$.

A lateral OPL gradient bends the transmitted ray. In the present software,
that effect is modeled with a thin-slab, weak-deflection approximation in
which the local angular deflection is proportional to the transverse
refractive-index gradient:
\begin{equation}
  \varepsilon_x
  \approx \frac{t}{n}\,\frac{\partial n}{\partial x},
  \label{eq:deflection}
\end{equation}
with the same expression used in the implemented scaling law. This model
assumes a sample of approximately uniform thickness, weak angular
deflection, and negligible variation through the optical depth.
This deflected ray appears to originate from a laterally shifted background
position. The apparent shift in the background plane is
\begin{equation}
  \Delta x_{\mathrm{bg}} = \varepsilon_x\,Z_d,
\end{equation}
and after projection through the imaging optics the corresponding
image-plane displacement in physical units is
\begin{equation}
  \Delta x_{\mathrm{img}} = \varepsilon_x\,Z_d\,\frac{z_i}{Z_B}.
  \label{eq:imgshift}
\end{equation}
In pixels, $u_{px} = \Delta x_{\mathrm{img}} / P_{px}$.

\paragraph{Decoupling limitation.}
Because BOS measures OPL, it cannot independently separate $n(x,y)$ from
$t(x,y)$ without additional information. The two reconstruction modes in
the software address this by assuming one quantity is known:

\begin{itemize}[noitemsep]
  \item \textbf{RI Variation (dn mode):} the sample thickness $t$ is assumed
        known and spatially uniform; the software reconstructs a relative
        refractive-index field $\Delta n(x,y)$.
  \item \textbf{Thickness Variation (mm mode):} the sample refractive index
        $n$ is assumed known and spatially uniform; the software reconstructs
        thickness variation $\Delta t(x,y)$.
\end{itemize}

Both modes are simplifications of the more general OPL measurement.
Any variation in the quantity assumed to be fixed will appear as a spurious
contribution to the reconstructed field.

\paragraph{Other physical contributions.}
Several effects beyond the OPL gradient influence the measured displacement
and should be considered when interpreting results:
\begin{itemize}[noitemsep]
  \item \textbf{Refraction at flat surfaces:} a ray entering a flat-faced
        sample at an angle undergoes lateral displacement by Snell's law
        even with no internal RI gradient. This geometric offset scales with
        sample thickness and off-axis angle, and is addressed by the optional
        field correction (Section~\ref{sec:refcorr}).
  \item \textbf{Sample tilt:} if the sample faces are not perpendicular to
        the optical axis, a systematic displacement offset is introduced
        across the entire field. This appears as a planar background in the
        $u$ and $v$ maps and in the reconstructed surface.
  \item \textbf{Defocus of background texture:} the effective background
        contrast changes slightly with RI variations, because the sample
        acts as a weak lens. In practice this effect is very small for weak
        disturbances.
\end{itemize}

\subsection{Displacement-to-Gradient Scaling}
\label{sec:corrscale}

The conversion from a raw pixel displacement to a physical gradient is
implemented in \cls{Session}\allowbreak\code{::ScaleFields()}. Combining
Equations~\eqref{eq:deflection} and~\eqref{eq:imgshift} gives
\begin{equation}
  u_{px} = \frac{\Delta x_{\mathrm{img}}}{P_{px}}
         = \frac{t}{n}\,\frac{\partial n}{\partial x}\,
           \frac{Z_d}{P_{px}}\frac{z_i}{Z_B},
\end{equation}
where, for readability in the equations below, $P_{px}$ is interpreted in
\si{mm/pixel}. In the UI and code implementation the pixel pitch is entered
in \si{\micro\metre/pixel}, so \cls{Session}\allowbreak\code{::ScaleFields()}
multiplies it by $10^{-3}$ before applying the same formulas.

Using
\begin{equation}
  \frac{z_i}{Z_B} = \frac{f}{Z_B-f},
\end{equation}
the RI-variation inversion becomes
\begin{equation}
  \frac{\partial n}{\partial x}
  =
  u_{px}\,
  P_{px}\,
  \frac{n\,(Z_B - f)}{f\,Z_d\,t}.
\end{equation}
The code therefore applies the correlation-map scale factor
\begin{equation}
  S_{\mathrm{corr}} =
    P_{px}\,\frac{n\,(Z_B - f)}{f\,Z_d\,t},
  \label{eq:scorr_dn}
\end{equation}
so that
\begin{equation}
  \frac{\partial n}{\partial x} = u_{px}\,S_{\mathrm{corr}}.
\end{equation}

For thickness-variation mode, the implemented model replaces $t$ with
$(n-1)$:
\begin{equation}
  S_{\mathrm{corr}} =
    P_{px}\,\frac{n\,(Z_B - f)}{f\,Z_d\,(n-1)}.
  \label{eq:scorr_mm}
\end{equation}

\subsection{Sample-Plane Grid Spacing}
\label{sec:gridspacing}

The displacement field is sampled on a correlation grid rather than on
individual pixels. The grid step (in image pixels) is
\begin{equation}
  \texttt{step} = \max\!\bigl(1,\; W - O\bigr),
  \label{eq:step}
\end{equation}
where $W$ is the window size and $O$ is the overlap. The corresponding
physical spacing in the sample plane is
\begin{equation}
  \Delta x_{\mathrm{grid}} =
    \texttt{step} \cdot P_{px} \cdot \frac{Z_a}{z_i}.
  \label{eq:dx}
\end{equation}

\subsection{Surface Scale Factor}
\label{sec:surfscale}

Because reconstruction integrates gradient values across discrete grid cells,
the reconstructed surface values must also absorb the physical cell width.
Therefore the surface scale factor is
\begin{equation}
  S_{\mathrm{surf}} = S_{\mathrm{corr}} \times \Delta x_{\mathrm{grid}}.
  \label{eq:ssurf}
\end{equation}
Both the displayed physical axis width and the reconstructed surface
amplitude therefore depend on the same grid-spacing factor. Correcting the
geometric parameters to fix the displayed axis also changes the reconstructed
surface amplitude.

\subsection{Optional Refraction Correction}
\label{sec:refcorr}

For samples with significant thickness, the curvature of off-axis rays
through the slab produces a background-plane displacement even for a
spectrally uniform sample. The software can compute and subtract this
geometric artifact via \cls{Session}\allowbreak\code{::ComputeRefractionCorrection()}.

For each correlator-grid point at column $c$, row $r$, the correction is
derived from Snell's law:
\begin{align}
  \theta_x &= \arctan\!\left(\frac{r_x}{Z_a}\right), &
  \theta_y &= \arctan\!\left(\frac{r_y}{Z_a}\right), \\
  \theta_{r,x} &= \arcsin\!\left(\frac{\sin\theta_x}{n}\right), &
  \theta_{r,y} &= \arcsin\!\left(\frac{\sin\theta_y}{n}\right),
\end{align}
where
\begin{equation}
  r_x = \left(c - \frac{w}{2}\right)\,\texttt{step}\,P_{px}\,\frac{Z_a}{z_i}
\end{equation}
is the sample-plane lateral position, with an analogous expression for
$r_y$.

The lateral displacement of the ray inside the slab is
\begin{equation}
  d_x = \frac{\sin(\theta_x - \theta_{r,x})}{\cos\theta_{r,x}}\,t,
\end{equation}
and similarly for $d_y$. Projecting back to the sensor plane and converting
to pixels gives the correction field that is subtracted before any further
scaling.

This correction should be considered an optional preprocessing step.
Its importance scales with both sample thickness and the departure of $n$
from unity: for samples with refractive indices close to 1, the Snell's-law offset is very small and
the correction has negligible impact on the final result. It becomes
progressively more important as $n$ increases or as $t$ grows large.

% ═══════════════════════════════════════════════════════════════════════════════
\section{Processing Pipeline}
\label{sec:pipeline}

Figure~\ref{fig:pipeline} shows the three-stage pipeline.

\begin{figure}[H]
  \centering
  \textit{(Pipeline diagram can be placed here.)}
  \caption{Overview of the Sand Dune processing pipeline.}
  \label{fig:pipeline}
\end{figure}

\subsection{Stage 1: Cross-Correlation}

\subsubsection{Image Preparation}

Images are loaded as normalized floating-point matrices by the \cls{Image}
class. Pixel values are mapped to $[0,1]$. The correlator operates directly
on these Eigen matrices without further preprocessing.

\subsubsection{Window Grid}

The image is divided into a regular grid of overlapping interrogation windows.
For a window size $W$ and overlap $O$, the grid step is given by
Equation~\eqref{eq:step}. The number of grid columns and rows is
\begin{equation}
  n_c = \left\lfloor \frac{W_{\mathrm{img}} - W}{\texttt{step}} \right\rfloor + 1,
  \qquad
  n_r = \left\lfloor \frac{H_{\mathrm{img}} - W}{\texttt{step}} \right\rfloor + 1.
\end{equation}

\subsubsection{FFT Cross-Correlation}

For each window pair (reference and flow), the correlator:
\begin{enumerate}[noitemsep]
  \item extracts the two interrogation blocks,
  \item skips the window if the local variance is below $10^{-3}$ (flat patch),
  \item subtracts the local mean from each window,
  \item computes the forward FFT of both windows,
  \item forms the cross-power spectrum
        $G = \overline{F_{\mathrm{ref}}}\,F_{\mathrm{flow}}$,
  \item applies the inverse FFT to obtain the raw correlation map,
  \item shifts the map so that zero displacement lies at the center.
\end{enumerate}

The resulting map $\phi$ satisfies
\begin{equation}
  \phi = \mathcal{F}^{-1}\!\left\{\overline{F_{\mathrm{ref}}}\,F_{\mathrm{flow}}\right\}.
\end{equation}

The signal-to-noise ratio is defined as
\begin{equation}
  \mathrm{s2n} = \frac{\phi_1}{\phi_2},
\end{equation}
where $\phi_1$ is the primary peak value and $\phi_2$ is the highest value
outside a $\pm 2$-pixel exclusion zone around the primary peak.

\subsubsection{CMM Sub-Pixel Refinement}
\label{sec:cmm}

Sub-pixel refinement is implemented using the Correlation Mapping Method
(CMM) of Chen and Katz~\cite{chen2005cmm}. The solver operates on a
$5 \times 5$ patch centred on the discrete peak.

The objective is to minimize the mean-normalized patch residual
\begin{equation}
  \varepsilon = \sum_{i=0}^{24}
    \left(\tilde{\phi}[i] - \tilde{\phi}'[i]\right)^2,
  \label{eq:cmm_obj}
\end{equation}
where $\tilde{\phi}[i] = \phi[i] / \bar{\phi}$ is the measured patch
normalized by its mean, and $\tilde{\phi}'[i]$ is the autocorrelation patch
evaluated at the sub-pixel offset $(u', v')$ via bicubic interpolation.
The autocorrelation patch is built once per window from the saved reference
FFT.

Minimization proceeds as follows:
\begin{enumerate}[noitemsep]
  \item \textbf{Seed:} a log-parabola fit along the dominant axis gives an
        initial $(u_0', v_0')$.
  \item \textbf{Gauss--Newton:} up to 8 iterations with backtracking line
        search (step sizes $\{1, 0.5, 0.25, 0.1, 0.05\}$). The Jacobian
        uses the analytic bicubic derivative.
  \item \textbf{Grid fallback:} if the Gauss--Newton solution does not improve
        the seed, a dense grid search on a $0.01$-pixel grid is used.
  \item \textbf{Boundary handoff:} if the solution is within $0.02$~px of
        $\pm 0.5$, the adjacent integer cell is tested and the lower-residual
        representation is kept.
\end{enumerate}

The solution is clamped to $\lvert u' \rvert \le 0.49$,
$\lvert v' \rvert \le 0.49$ to prevent ambiguous half-pixel values.

\paragraph{Peak locking.}
A key motivation for using CMM rather than a simpler centroid or parabolic
fit is the suppression of peak locking. Conventional sub-pixel estimators
that fit a symmetric function to the discrete correlation peak are biased
toward integer and half-integer values, producing characteristic spikes in
the displacement histogram. CMM avoids this by matching the entire local
patch shape to a reference autocorrelation rather than fitting only the peak
tip. The boundary-handoff logic and the clamp at $\pm 0.49$~px are
additional safeguards that prevent the estimate from locking to cell
boundaries. The peak-locking fraction is tracked explicitly in the test
suite (Section~\ref{sec:validation}) to confirm that this behaviour
remains controlled.

\begin{figure}[H]
  \centering
  \figimg[height=12cm]{peak_locking.png} % replace with actual filename
  \caption{An example of peak locking patterns formed from a soft gradient measurement field.}
  \label{fig:panel_load}
\end{figure}

\subsubsection{Progressive Image Deformation (Optional)}
\label{sec:pid}

When \param{enable\_pid} is active, the correlator performs Progressive
Image Deformation (PID) to account for local shear and strain in the
displacement field. PID iteratively warps the interrogation windows to
match the local deformation model, so that each correlation pass measures
only the residual displacement after the current best estimate is removed.

\paragraph{Algorithm.} Starting from the baseline rigid-window correlation:
\begin{enumerate}[noitemsep]
  \item Compute first-order displacement gradients
        $\partial u/\partial x$, $\partial u/\partial y$,
        $\partial v/\partial x$, $\partial v/\partial y$
        via central differences at the correlator-grid spacing.
  \item Apply \param{pid\_smoothing\_passes} passes of a $3{\times}3$ box
        filter to the deformation field.
  \item For each window, construct an affine warp:
        $\mathbf{x}_{\mathrm{warped}} = \mathbf{x} + \mathbf{u}(\mathbf{x}_c)
        + J\,(\mathbf{x} - \mathbf{x}_c)$,
        where $\mathbf{x}_c$ is the window centre and $J$ is the local
        deformation Jacobian.
  \item Resample the flow window using bilinear interpolation and
        re-correlate to measure the residual displacement.
  \item Accept the update for a window only if the CMM residual improves
        (or the S2N improves when residuals are not available).
  \item Repeat for \param{pid\_iterations} additional passes. If the
        mean correction across all windows falls below $10^{-3}$~px,
        terminate early.
\end{enumerate}

The \param{pid\_relaxation} factor damps the accepted residual updates to
improve convergence stability.

\subsection{Stage 2: Validation}
\label{sec:validation_stage}

Validation is implemented in \file{src/grains/validation.cpp} and uses a
two-pass outlier-detection strategy.

\paragraph{Pass 1: Signal-to-noise threshold.}
A vector is marked invalid if
\begin{equation}
  \mathrm{s2n} < 1.3.
\end{equation}

\paragraph{Pass 2: Normalized residual (UNO criterion).}
For each vector the $8$-connected neighbourhood median is computed:
\begin{equation}
  u_{\mathrm{med}} = \operatorname{median}\!\bigl(u_k,\; k \in \mathcal{N}_8\bigr).
\end{equation}
The normalized deviation is
\begin{equation}
  \tilde{u}_{\mathrm{nrm}} =
    \frac{|u - u_{\mathrm{med}}|}
         {\operatorname{median}(|u_k - u_{\mathrm{med}}|) + \varepsilon},
  \qquad \varepsilon = 0.1,
\end{equation}
and the vector is flagged if
\begin{equation}
  \sqrt{\tilde{u}_{\mathrm{nrm}}^2 + \tilde{v}_{\mathrm{nrm}}^2} > 2.0.
\end{equation}
This is the universal outlier detection (UNO) criterion of Westerweel and
Scarano~\cite{westerweel2005uno}.

\paragraph{Replacement.} All flagged vectors are replaced with the
$8$-neighbourhood median of the unflagged neighbours.

\subsection{Stage 3: Surface Reconstruction}
\label{sec:reconstruction}

\subsubsection{Masked Sparse Poisson Solver (Default Method)}

The default reconstruction method builds and solves a sparse linear system
on the masked correlation grid. Let $s_{i,j}$ be the unknown scalar field on
valid nodes. Edge gradients are formed as
\begin{align}
  g^x_{i,j} &= \tfrac{1}{2}\!\left(u_{i,j} + u_{i,j+1}\right), \\
  g^y_{i,j} &= \tfrac{1}{2}\!\left(v_{i,j} + v_{i+1,j}\right),
\end{align}
and the energy to be minimized is
\begin{equation}
  E(s) =
  \sum_{\text{x-edges}}
    \!\left(s_{i,j+1} - s_{i,j} - g^x_{i,j}\right)^2
  +
  \sum_{\text{y-edges}}
    \!\left(s_{i+1,j} - s_{i,j} - g^y_{i,j}\right)^2.
  \label{eq:poisson_energy}
\end{equation}

Setting $\partial E/\partial s_{i,j} = 0$ yields a standard sparse
Poisson-style linear system. The domain is restricted to nodes inside the
mask, so no artificial boundary data is inserted. The system is assembled
with Eigen triplets and solved using
\allowbreak\code{Eigen::ConjugateGradient} with tolerance $10^{-5}$ and a maximum of
$\max(4N,\,200)$ iterations, where $N$ is the number of unknowns.

The constant-offset nullspace is resolved by a soft gauge condition on the
first valid interior node (pinned to zero). The reconstructed field is
therefore relative in absolute value but meaningful in shape.

\subsubsection{Frankot--Chellappa (Optional Method)}

The software also provides an FFT-based reconstruction option following
Frankot and Chellappa~\cite{frankot1988}. This path is exposed in the UI as
an alternative integrator and is implemented by \allowbreak\code{ComputeFC()}.
It mirrors the displacement field into a
$2{\times}2$-tiled domain to enforce periodic boundary conditions and then
integrates in the frequency domain:
\begin{equation}
  \hat{s}(f_x, f_y) =
    \frac{i f_x\,\hat{u}(f_x,f_y) + i f_y\,\hat{v}(f_x,f_y)}
         {2\pi(f_x^2 + f_y^2) + \varepsilon}.
\end{equation}

On full-domain inputs, the Frankot--Chellappa method can produce results
comparable to the sparse Poisson solver and is substantially faster, since
it requires only a small number of FFTs regardless of grid size. Its main
caveat in the present BOS workflow is the treatment of the sample mask.
The current implementation applies the mask by zeroing $u$ and $v$ outside
the valid region before the FFT integration. This is convenient and keeps
the option usable in the UI, but it does not treat the exterior as an
unknown domain in the same way as the sparse solver. As a result, sharp mask
boundaries can introduce artificial low-frequency structure near the sample
edge. For masked aerogel-style datasets, the sparse Poisson solver is
therefore the default and generally the more trustworthy choice, while
Frankot--Chellappa remains useful as a fast comparison method on cleaner or
less heavily masked fields.

% ═══════════════════════════════════════════════════════════════════════════════
\section{Software Architecture}
\label{sec:software}

\subsection{Build System and Dependencies}

The project is built with CMake (C++23 required). External dependencies are:

\begin{center}
\begin{tabular}{ll}
  \toprule
  \rowcolor{accentlight} Library & Purpose \\
  \midrule
  SDL3, SDL3\_image & Application window, rendering, image file I/O \\
  Dear ImGui        & All UI panels and controls \\
  ImPlot            & Heatmap visualizations \\
  Eigen3            & Dense/sparse linear algebra \\
  FFTW3 (float)     & FFT cross-correlation and Frankot--Chellappa \\
  doctest           & Unit-test framework (test binary only) \\
  \bottomrule
\end{tabular}
\end{center}

The build produces two targets:
\begin{itemize}[noitemsep]
  \item \allowbreak\code{sand\_dune}: the GUI application,
  \item \allowbreak\code{run\_tests}: the standalone test binary.
\end{itemize}

\subsection{Module Organization}

\begin{center}
\begin{tabular}{lll}
  \toprule
  \rowcolor{accentlight} Namespace / directory & Classes & Role \\
  \midrule
  \file{src/main.cpp}         &              & Application entry point \\
  \file{include/session.h},   & \cls{Session}& Pipeline orchestration, scaling, \\
  \file{src/session.cpp}      &              & async execution, CSV export \\
  \file{include/grains/}      & \cls{Correlator} & FFT correlation, CMM, PID \\
  \file{src/grains/}          & \cls{Validation} & Outlier detection \& replacement \\
                              & \cls{Reconstruction} & Poisson / Frankot--Chellappa integration \\
  \file{include/water/}       & \cls{UI}, \cls{App} & ImGui panels, SDL3 shell \\
  \file{src/water/}           &              & \\
  \file{include/sun/}         & \cls{Image}, \cls{Mask} & Image loading, mask generation \\
  \file{src/sun/}             &              & \\
  \file{include/wind/}        & \cls{VectorField} & Displacement-field container \\
  \file{src/wind/}            & (parameter structs) & \\
  \file{tests/}               &              & Development verification tests \\
  \bottomrule
\end{tabular}
\end{center}

\subsection{Session: Pipeline Orchestrator}

\cls{Session} (\file{include/session.h}) is the top-level controller.
It owns the input images, all intermediate data arrays, the parameter
structs, and the stage-state machine:

\begin{cppcode}
// Images
Image ref;
std::vector<Image> flows;

// Intermediate data
std::vector<VectorField> raw_correlation_field;
std::vector<VectorField> correlation_field;
std::vector<VectorField> raw_val_field;
std::vector<VectorField> val_field;
std::vector<Eigen::MatrixXf> raw_surface;
std::vector<Eigen::MatrixXf> surface;
Eigen::MatrixXf correction[2];  // refraction correction (u, v)

// Parameters
CorrelatorParameters correlatorparameters;
OpticalParameters    opticalparameters;
Mask                 mask;

// Stage state (per stage: Idle, Ready, Busy, Done, Dirty)
std::atomic<StageState> stagestates[STAGE_TOTAL];
\end{cppcode}

The three pipeline stages run asynchronously via \allowbreak\code{std::async}.
The UI polls the stage states to display progress and enables/disables
run buttons accordingly.

\subsection{Correlator}

\cls{Correlator} (\file{include/grains/correlator.h}) takes a reference
\cls{Image} and a flow \cls{Image} and returns a \cls{VectorField}.
Internally it allocates FFTW plans once at construction time and reuses
the same transform buffers for every window.

\begin{cppcode}
struct CorrelatorParameters {
    int   window_size        = 32;
    int   overlap            = 24;
    bool  enable_pid         = false;
    int   pid_iterations     = 2;
    float pid_relaxation     = 1.0f;
    int   pid_smoothing_passes = 1;
};
\end{cppcode}

The main entry point is \allowbreak\code{Compute(ref, flow)}, which dispatches to
\allowbreak\code{ComputeWithPid()} or \allowbreak\code{Compute\-SinglePass()} based on \param{enable\_pid}.

\subsection{Reconstruction}

\cls{Reconstruction} (\file{include/grains/reconstruction.h}) provides two
integration back ends. The default path takes a \cls{VectorField} and a
binary mask \allowbreak\cls{Eigen\-::MatrixXi}, builds an index map that assigns
consecutive unknowns only to valid interior cells, constructs the sparse
system from edge-gradient equations, and calls the Eigen conjugate-gradient
solver. An FFT-based Frankot--Chellappa path is also available as an
alternative reconstruction option.

\subsection{VectorField}

\cls{VectorField} (\file{include/wind/vectorfield.h}) stores three
\allowbreak\code{Eigen::MatrixXf} matrices ($u$, $v$, s2n) and grid dimensions.
It exposes a \allowbreak\code{SaveCSV(path)} method that writes columns \texttt{col},
\texttt{row}, \texttt{u}, \texttt{v}, \texttt{s2n} in column-major order.

% ═══════════════════════════════════════════════════════════════════════════════
\section{Using the Application}
\label{sec:usage}

\subsection{Typical Workflow}

\begin{enumerate}
  \item \textbf{Load images.} Use the \textbf{Load} panel to browse for a
        reference image and one or more disturbed (flow) images.
        All images must share the same pixel dimensions.

  \item \textbf{Set parameters.} In the \textbf{Parameters} panel, enter
        the optical geometry ($Z_d$, $Z_a$, $f$, $P_{px}$), sample
        properties ($n$, $t$), and the correlation settings (window size,
        overlap, PID toggle). The \textbf{Calculations} panel updates in
        real time to show the derived grid spacing, physical field width,
        and estimated feature size; use these to check that the window
        size is sensible for the expected background texture.

  \item \textbf{Configure the mask.} In the \textbf{Parameters} panel,
        enable the circular mask and set its centre and radius to enclose
        the sample region. Reconstruction will be confined to this domain.

  \item \textbf{Run the pipeline.} In the \textbf{Pipeline} panel, click
        \textbf{Run Full Pipeline} to run all three stages sequentially, or run
        each stage individually. A progress bar and ETA are shown for
        each stage. With multiple flow images, all frames are processed
        in the same batch.

  \item \textbf{Inspect results.} Switch between the \textbf{Correlation},
        \textbf{Validation}, and \textbf{Surface} views using the tabs in
        the visualization panel. For correlation and validation, select
        the $u$, $v$, or s2n map, and toggle between raw pixel units and
        scaled physical units. Adjust the colormap range using the
        min/max sliders.

  \item \textbf{Export.} In the \textbf{Save} panel, choose an output
        directory and click \textbf{Save}. Three CSV files are written:
        \texttt{\_correlation.csv}, \texttt{\_val.csv}, and
        \texttt{\_surface.csv}. For multi-frame batches, a numeric suffix
        is appended per frame.
\end{enumerate}

\subsection{Panel Reference}

\subsubsection{Load Panel}

Provides file-picker buttons to select the reference image and one or
more flow images. The flow image list supports multi-frame batches.
When multiple flow images are loaded, navigation arrows are shown in
the visualization panels to step through frames.

\subsubsection{Parameters Panel}

\begin{center}
\begin{tabular}{lll}
  \toprule
  \rowcolor{accentlight} Group & Control & Effect \\
  \midrule
  Correlation & Window size & Interrogation window side length (pixels) \\
              & Overlap     & Window overlap (pixels; step = size $-$ overlap) \\
              & Enable PID  & Enables iterative window deformation \\
              & PID iterations & Number of PID passes after baseline \\
              & PID relaxation & Damping on residual updates ($[0,1]$) \\
              & PID smoothing  & $3{\times}3$ smoothing passes on deformation field \\
  \midrule
  Optical     & $P_{px}$       & Sensor pixel pitch (\si{\micro\metre}) \\
              & $Z_d$          & Background-to-sample distance (\si{mm}) \\
              & $Z_a$          & Sample-to-lens distance (\si{mm}) \\
              & $f$            & Lens focal length (\si{mm}) \\
  \midrule
  Sample      & $n$            & Sample refractive index \\
              & $t$            & Sample thickness (\si{mm}) \\
  \midrule
  Reconstruction & Mode        & RI variation (dn) or thickness variation (mm) \\
                 & Apply correction & Subtract Snell's-law refraction offset \\
  \midrule
  Mask        & Enable mask    & Restrict reconstruction to circular ROI \\
              & Centre (x, y)  & Mask centre in correlation-grid coordinates \\
              & Radius         & Mask radius in correlation-grid cells \\
  \bottomrule
\end{tabular}
\end{center}


\subsubsection{Calculations Panel}

Displays derived quantities updated in real time from the current
parameters: correlation-grid step (px), sample-plane grid spacing (mm),
physical field width (mm), image magnification, and the expected
pixel displacement for a reference gradient magnitude.


\subsubsection{Pipeline Panel}

Shows three stage rows (Correlation, Validation, Reconstruction) each with
a \textbf{Run} button, a progress bar, and a status label (Idle / Ready /
Busy / Done / Dirty). A \textbf{Run Full Pipeline} button chains all three stages.

\subsubsection{Visualization Panel}

\begin{itemize}[noitemsep]
  \item \textbf{Correlation / Validation tabs:} heatmap of the $u$, $v$,
        or s2n map. A toggle switches between raw pixels and scaled
        physical units. Min/max sliders control the colormap range.
  \item \textbf{Surface tab:} heatmap of the reconstructed scalar field.
        Axes are labelled in physical millimetres. Min/max sliders control
        the colormap range. A tooltip shows the value under the cursor.
\end{itemize}

\subsubsection{Save Panel}

Allows choosing an output directory (via system file picker) and exports
the current results to CSV. Saving runs asynchronously to avoid blocking
the UI.

\subsubsection{Settings Panel}

Offers a dark/light theme toggle.

\subsection{CSV File Formats}

\paragraph{Correlation and Validation CSV.}
The correlation and validation result files share the following layout:
\begin{fmtcode}
rows, cols
<height>, <width>
u, v, s2n
<u00>, <v00>, <s2n00>
...
\end{fmtcode}
Values are in physical units when optical parameters have been set;
otherwise they are in raw pixels.

\paragraph{Surface CSV.}
The surface file uses column-major ordering:
\begin{fmtcode}
rows, cols
<height>, <width>
<z[0,0]>
<z[1,0]>
...
\end{fmtcode}
Values are stored in column-major order (all rows of column 0, then
all rows of column 1, etc.) in the units determined by the selected
reconstruction mode (dn or mm).

% ═══════════════════════════════════════════════════════════════════════════════
\section{Validation, Precision, and Repeatability}
\label{sec:validation}

\subsection{Overview}

Validation is performed at two levels. At the vector level, the pipeline
uses signal-to-noise thresholding and neighbourhood-median outlier detection
(Section~\ref{sec:validation_stage}) to clean each displacement field in-line.
At the software level, the development verification suite in
\file{tests/test.cpp} was used to characterize numerical behaviour,
estimate precision, and check for artifacts during implementation. The test
binary (\allowbreak\code{run\_tests}) is not part of the normal user
workflow; it exists as a development and regression tool.

\subsection{Image-Based Precision Metric}

The most important quantitative precision figure available for the current
system comes from the measured image pair \texttt{if\_0.1}, stored under
\file{images/slides}. This image pair was taken of a glass optical window 
of known uniform refractive index and thickness. The expected
displacement field is smooth and nearly planar, so the residual from a
best-fit plane directly quantifies spurious structure such as striping,
local roughness, and sub-pixel locking artifacts.

During software verification, the test \textit{Correlation Repo Image Pair}
fits a least-squares plane to the recovered $u$ and $v$ fields and reports
the RMS deviation. Results from the current codebase are summarised below.

\begin{center}
\begin{tabular}{lll}
  \toprule
  \rowcolor{accentlight} Metric & $u$ (px) & $v$ (px) \\
  \midrule
  Whole-map RMS, no PID  & 0.659\,215 & 0.669\,790 \\
  Whole-map RMS, PID     & 0.595\,992 & 0.579\,117 \\
  Central-crop RMS, no PID & 0.096\,441 & 0.105\,765 \\
  Central-crop RMS, PID    & 0.048\,692 & 0.045\,640 \\
  \bottomrule
\end{tabular}
\end{center}

The central crop removes the outer 10\,\% border on each side (containing 
ringing and artifacts from the sample holder and image edge), leaving the
central 80\,\% of the map, which is the most representative region for the
present system. Enabling PID reduces the central-crop RMS by approximately
\textbf{49.5\,\%} in $u$ and \textbf{56.8\,\%} in $v$.

\subsection{Artifact Diagnostics}

The development verification suite also includes image-based diagnostics
designed to identify structured artifacts in smooth gradient fields:
\begin{itemize}[noitemsep]
  \item \textbf{Striping check:} the dominant-axis autocorrelation of the
        plane-fit residual field is required to remain below a threshold
        ($< 0.75$), and the orthogonal-axis correlation must be low
        ($< 0.10$), to confirm that no spatial periodicity is present.
  \item \textbf{Peak-locking check:} the fraction of recovered displacement
        values within $0.03$~px of an integer or half-integer is measured.
        Excessive concentration at these values indicates sub-pixel bias.
  \item \textbf{Validated residual ratio:} the ratio of the validation-stage
        residual to the correlation residual is tracked to confirm that
        the validation step is not masking systematic errors.
\end{itemize}

\subsection{Supplementary Software Checks}

In addition to the measured-image verification above, the
\textit{Correlator Computation} test case in \file{tests/test.cpp} was used
during development to verify expected behaviour on controlled synthetic
inputs. These checks are useful regression guards, but they are not the
primary experimental precision metric for the present system.

\paragraph{Zero displacement.}
A random-noise image is correlated against itself. The measured RMSE
against the zero-displacement ground truth is $0.0000$~px, with a
maximum absolute displacement of $0.0000$~px anywhere on the grid.
This confirms that the correlator introduces no systematic bias on a
flat-spectrum input.

\paragraph{Sub-pixel translation.}
A synthetic spot-field (120 Gaussian spots, $\sigma = 1.1$~px, rendered
on a $192 \times 192$ grid) is shifted by bilinear interpolation for four
imposed offsets and the per-shift RMSE against the known ground-truth
displacement is reported:

\begin{center}
\begin{tabular}{ll}
  \toprule
  \rowcolor{accentlight} Imposed shift (px) & RMSE (px) \\
  \midrule
  $(+0.25,\ -0.30)$ & 0.0633 \\
  $(+0.10,\ +0.12)$ & 0.0254 \\
  $(-0.35,\ +0.28)$ & 0.0711 \\
  $(+0.42,\ -0.18)$ & 0.0714 \\
  \bottomrule
\end{tabular}
\end{center}

All four cases remain well below the $0.08$~px acceptance threshold,
confirming that the CMM sub-pixel solver operates correctly across the
$\pm 0.5$~px cell.

\paragraph{PID affine deformation.}
A spot-field (160 spots, same parameters) is warped by an affine
transformation with both translation and first-order shear:
$u_0 = 1.15$~px, $v_0 = -0.75$~px, $\partial u/\partial x = 0.0105$,
$\partial u/\partial y = 0.0040$, $\partial v/\partial x = -0.0035$,
$\partial v/\partial y = 0.0085$. The RMSE against the analytical
displacement field is compared across iteration counts:

\begin{center}
\begin{tabular}{ll}
  \toprule
  \rowcolor{accentlight} Method & RMSE (px) \\
  \midrule
  Baseline (no PID) & 0.2745 \\
  PID (2 iterations) & 0.0875 \\
  PID (3 iterations) & 0.0846 \\
  \bottomrule
\end{tabular}
\end{center}

Two PID iterations reduce the RMSE by $68.1\,\%$, from $0.2745\,\mathrm{px}$
to $0.0875\,\mathrm{px}$. The residual is sub-pixel and consistent with the
correlation noise floor, indicating that the deformation correction removes
systematic error rather than merely redistributing it. A third iteration
yields only a further $3.3\,\%$ reduction ($0.0846\,\mathrm{px}$),
demonstrating the expected diminishing returns: most of the correctable
deformation error is captured in the first two passes.

\subsection{Repeatability}

Repeatability is evaluated experimentally by acquiring repeated image sets
under different distance and rotation conditions to asses the measured field:
\begin{itemize}[noitemsep]
  \item \textbf{Rotation symmetry:} the repeated acquisition of the sample was 
  done slightly rotated to assess overall gradient symmetry and mapping quality.
  \item \textbf{Distance symmetry:} acquisitions at controlled $Z_A$ distances were done to 
  assess optical parameter impact and repeatability between such changes.
\end{itemize}

\begin{figure}[H]
  \centering
  \begin{minipage}[t]{0.48\linewidth}
    \centering
    \figimg[width=\linewidth]{repeatability_far.png}
    \caption{Refractive index map for a further placement from sample to lens.}
    \label{fig:rep_far}
  \end{minipage}\hfill
  \begin{minipage}[t]{0.48\linewidth}
    \centering
    \figimg[width=\linewidth]{repeatability_near.png}
    \caption{Refractive index map for a near placement of the same sample.}
    \label{fig:rep_near}
  \end{minipage}
\end{figure}


\begin{figure}[H]
  \centering
  \begin{minipage}[t]{0.48\linewidth}
    \centering
    \figimg[width=\linewidth]{repeatability_far_cor.png}
    \caption{Refractive index gradient map for a further placement from sample to lens.}
    \label{fig:rep_rotation}
  \end{minipage}\hfill
  \begin{minipage}[t]{0.48\linewidth}
    \centering
    \figimg[width=\linewidth]{repeatability_near_cor.png}
    \caption{Refractive index gradient map for a near placement of the same sample.}
    \label{fig:rep_distance}
  \end{minipage}
\end{figure}

The four figures above showed good repeatability and consistency between both changing 
setup geometries and reference image changes, as each used difference reference patterns for correlation and measurement.

% ═══════════════════════════════════════════════════════════════════════════════
\section{System Precision Estimate}
\label{sec:precision}

This section quantifies the minimum resolvable refractive-index gradient for
the current optical configuration and places it in the context of the weakest
gradients expected in practice.

\subsection{Current Optical Setup}

The figures below are computed for the following configuration:

\begin{center}
\begin{tabular}{lll}
  \toprule
  \rowcolor{accentlight} \textbf{Parameter} & \textbf{Symbol} & \textbf{Value} \\
  \midrule
  Lens focal length        & $f$       & \SI{25}{mm} \\
  Background--sample dist. & $Z_d$     & \SI{220}{mm} \\
  Sample--lens distance    & $Z_a$     & \SI{45}{mm} \\
  Total object distance    & $Z_B$     & \SI{265}{mm} \\
  Pixel pitch              & $P_{px}$  & \SI{2.315}{\micro\metre} \\
  Sample refractive index  & $n$       & $1.08$ \\
  Sample thickness         & $t$       & \SI{0.924}{mm} \\
  \bottomrule
\end{tabular}
\end{center}

\subsection{Minimum Resolvable Gradient}

The minimum detectable displacement is taken as three times the PID
central-crop plane-fit RMS from the slide validation tests
(Section~\ref{sec:validation}):
\begin{equation}
  u_{\min} = 3 \times 0.049 \approx 0.15~\text{px}.
\end{equation}
Using the three-sigma criterion gives a conservative threshold that
distinguishes genuine signal from measurement noise at high confidence.

The displacement-to-gradient scale factor (Equation~\eqref{eq:scorr_dn}),
evaluated in SI units throughout, is
\begin{equation}
  S_{\mathrm{corr}}
  = \frac{P_{px}\,(Z_B - f)\,n}{f\,Z_d\,t}
  = \frac{(2.315\times10^{-6})(240\times10^{-3})(1.08)}
         {(25\times10^{-3})(220\times10^{-3})(0.924\times10^{-3})}
  \approx 0.1181~\mathrm{m}^{-1}\,\mathrm{px}^{-1}.
\end{equation}

The minimum resolvable refractive-index gradient is therefore
\begin{equation}
  \left.\frac{dn}{dx}\right|_{\min}
  = S_{\mathrm{corr}} \times u_{\min}
  = 0.1181 \times 0.15
  \approx 1.77\times10^{-2}~\mathrm{m}^{-1}
  = 1.77\times10^{-5}~\mathrm{dn/mm}.
\end{equation}

\subsection{Expected Worst-Case Gradient}

As a benchmark, consider a transparent sample with a total refractive-index
variation of $\Delta n = 0.005$ distributed uniformly across the full sample
diameter of \SI{24}{mm}. This represents an internally measured contrast
value assumed to be spread as a slow, gradual gradient over the entire
aperture, which constitutes the worst case in the sense that the gradient
is as shallow as it can be for a given total RI excursion:
\begin{equation}
  \left.\frac{dn}{dx}\right|_{\mathrm{wc}}
  = \frac{\Delta n}{d}
  = \frac{0.005}{24}
  \approx 2.08\times10^{-4}~\mathrm{dn/mm}.
\end{equation}

\subsection{Detection Margin}

The ratio of the expected worst-case gradient to the minimum resolvable
gradient gives the system margin:
\begin{equation}
  \text{margin}
  = \frac{(dn/dx)_{\mathrm{wc}}}{(dn/dx)_{\min}}
  = \frac{2.08\times10^{-4}}{1.77\times10^{-5}}
  \approx 11.8.
\end{equation}

Even in the worst case of a maximally shallow gradient, the current
system can resolve the expected RI variation by approximately one order
of magnitude. Steeper local gradients (more concentrated RI variations)
would be resolved with even greater margin. This confirms that the
optical sensitivity of the current setup is well matched to the target
measurement regime.

% ═══════════════════════════════════════════════════════════════════════════════
\section{Calibration and Interpretation}
\label{sec:calibration}

The most sensitive practical issue is geometric calibration, specifically
setting $Z_a$ correctly. In the software model, $Z_a$ is an effective optical
distance, not simply a mechanical measurement from the lens barrel. An
incorrect $Z_a$ shifts both the displayed field width and the reconstructed
surface amplitude simultaneously, because both depend on the same
grid-spacing factor $\Delta x_{\mathrm{grid}}$ (Equation~\eqref{eq:dx}).

The recommended calibration procedure is therefore not a direct ruler
measurement. Instead:
\begin{enumerate}[noitemsep]
  \item Place a target of known physical width in the background plane.
  \item Record how many pixels span that width.
  \item Infer the effective magnification and back-calculate $Z_a$ so that
        the reported field width in the \textbf{Calculations} panel matches
        the known dimension.
\end{enumerate}

Once $Z_a$ is correctly calibrated, the surface amplitude will also be
on the correct absolute scale.

The reconstructed surface is always relative in absolute value (the gauge
constant is arbitrary), so absolute-offset comparisons between separate
experiments require an independent reference point.

\subsection{Interpreting Results: Reconstruction vs.\ Gradient Maps}

The reconstruction and the raw displacement maps each reveal different
aspects of the optical quality of a transparent sample, and both should
be examined together rather than relying on either alone.

The reconstructed surface (dn or thickness map) gives a global view of
the spatial distribution of optical path length across the sample.
Broad, smooth features such as overall refractive-index gradients,
wedge-shaped thickness variation, or large-scale density non-uniformities
appear clearly in the integrated field. This view is well suited to
comparing samples, measuring the amplitude of a known disturbance, or
checking that the reconstructed field has the expected large-scale shape.

The correlation $u$ and $v$ gradient maps, by contrast, directly show
where harsh or rapid variations in refractive index occur within the
sample. Sharp features such as inclusions, bubbles, cracks, or regions
of steep RI gradient produce strong, localised displacement peaks in these
maps that may be smeared or suppressed by the spatial-integration step of
the reconstruction. Examining the gradient maps in addition to the surface
is therefore important for identifying optically defective regions that
a smooth reconstruction might underrepresent.

In practice: use the reconstructed surface to characterise the overall
optical uniformity and quantify the amplitude of large-scale variations;
use the $u$ and $v$ maps to locate and assess localised optical defects.

% ═══════════════════════════════════════════════════════════════════════════════
\section{Known Caveats}
\label{sec:caveats}

\subsection{Fundamental Measurement Limitations}

\begin{itemize}

  \item \textbf{Refractive index and thickness cannot be decoupled.}
        BOS measures the integrated optical path length (OPL) through the
        sample. A thicker region with uniform refractive index produces
        exactly the same deflection signal as a thinner region with a
        higher refractive index. Without an independent thickness
        measurement (e.g.\ profilometry, OCT), the two contributions
        cannot be separated from the BOS data alone.

  \item \textbf{The reconstructed map is not a pure refractive-index map.}
        Although the output is labelled in refractive-index units via the
        Gladstone-Dale scaling, the deflection signal contains contributions
        from all sources that bend the light ray:
        \begin{itemize}[noitemsep]
          \item refractive index variation within the sample bulk,
          \item physical tilt or wedge of the sample surfaces (a flat but
                tilted window deflects light even with perfectly uniform RI),
          \item surface curvature and figure errors,
          \item internal stress birefringence (though BOS is not
                polarisation-resolved and therefore only captures the
                scalar mean),
          \item any other optical inhomogeneity along the line of sight.
        \end{itemize}
        All of these contribute to the measured gradient field and are
        integrated into the final surface by the reconstruction step. There
        is no way to attribute a given feature in the output map to one
        specific physical cause without additional independent measurements.

  \item \textbf{BOS is a line-of-sight-integrated measurement.}
        The technique collapses the full 3D refractive-index field onto a
        single 2D projection. Any variation along the optical axis (depth
        within the sample) is summed and cannot be resolved. Tomographic
        BOS can recover 3D structure but requires multiple viewing angles
        and is not implemented here.

  \item \textbf{Only in-plane deflections are measured.}
        A refractive-index gradient along the optical axis ($\partial n /
        \partial z$, i.e.\ along the camera axis) does not deflect the ray
        laterally and therefore produces no signal. BOS is sensitive only
        to gradients in the plane perpendicular to the optical axis.

  \item \textbf{The measurement is a gradient, not a direct surface.}
        The correlation pipeline recovers $\partial(\Delta n \cdot t) /
        \partial x$ and $\partial(\Delta n \cdot t) / \partial y$. The
        surface map is obtained by integrating these gradients. Integration
        accumulates errors: a spatially correlated noise floor in the
        gradient field can produce low-frequency artefacts in the
        reconstructed surface that are not present in the gradient maps.
        Always inspect both the gradient fields and the reconstruction
        when diagnosing anomalies.

  \item \textbf{Small-angle (paraxial) approximation.}
        The deflection-angle formula $\varepsilon \approx \tan\varepsilon$
        is used throughout. For typical BOS setups with small $Z_d / Z_a$
        ratios and moderate sample-induced deflections, the error is
        negligible. For large deflections or extreme geometry, the
        approximation breaks down and a full ray-tracing model is required.

  \item \textbf{The reference image is assumed deflection-free.}
        Any structured distortion present in the reference image
        (lens distortion, ambient air turbulence, vibration during
        acquisition) will appear as a systematic offset in every result
        derived from that reference. Lens distortion should be corrected
        or characterised separately; the reference image should be acquired
        under the same thermal and vibration conditions as the sample
        images.

  \item \textbf{The gradient field is not guaranteed to be curl-free.}
        A physically consistent surface has $\partial u / \partial y =
        \partial v / \partial x$ everywhere. Measurement noise, outliers,
        and genuine optical effects (e.g.\ diffraction at edges) produce a
        curl component. The sparse Poisson solver discards this component
        implicitly by finding the least-squares integrable surface. The
        discarded curl energy is not reported, so a high-curl field can
        reconstruct to a smooth-looking surface while hiding significant
        measurement inconsistency.

\end{itemize}

\subsection{Software and Implementation Caveats}

\begin{itemize}

  \item \textbf{Frankot--Chellappa remains an option, but its mask treatment
        is weaker.}
        The default reconstruction path is the masked sparse Poisson solver.
        Frankot--Chellappa can still be selected in the UI, but in the
        current implementation the mask is enforced by zeroing the gradient
        field outside the sample before the FFT solve. This can introduce
        boundary-driven artifacts that the sparse masked-domain solver avoids.

  \item \textbf{Reconstructed surface is relative.}
        The constant-offset degree of freedom is resolved by a gauge
        condition on one interior node. The absolute offset has no
        physical meaning.

  \item \textbf{Field correction is applied before the ``raw'' correlation
        view.}
        The stored raw correlation field already has the Snell's-law
        correction subtracted when that option is enabled. The ``raw''
        label in the UI refers to pixel-unit (unscaled) data, not
        uncorrected data.

  \item \textbf{Validation criteria are not uniformly applied across all
        code paths.}
        The main post-processing path uses both S2N and UNO thresholding.
        Some alternative overloads apply only the S2N pass.

  \item \textbf{Reconstruction progress granularity.}
        Each full sparse solve is counted as one progress step. For small
        masks or small grids, the progress bar may jump rather than
        advance smoothly.

  \item \textbf{Geometry calibration is critical.}
        An incorrectly set $Z_a$ shifts both the spatial axis and the
        surface amplitude by the same factor. Always calibrate against a
        known physical target before interpreting absolute values.

\end{itemize}

% ═══════════════════════════════════════════════════════════════════════════════
\section{Conclusion}

Sand Dune is a complete BOS measurement system for mapping weak optical
non-uniformities in transparent media at high spatial resolution.
Background pattern design, optical layout, calibration, cross-correlation,
validation, and reconstruction all contribute directly to the quality of the
final result. The current software combines FFT-based
windowed cross-correlation, CMM sub-pixel refinement (with optional PID
window deformation), UNO-based vector validation, and a selectable
reconstruction stage with masked sparse Poisson integration as the default
and Frankot--Chellappa as an alternative option.

At present, the most meaningful precision figure is the best-fit-plane RMS
obtained from the measured \texttt{if\_0.1} dataset on the central crop:
approximately $0.049$~px in $u$ and $0.046$~px in $v$ with PID enabled.
These figures provide a quantitative development-verification benchmark for
the current implementation and establish the noise floor used elsewhere in
this document. The next additions to the report should be the final
repeatability figures and the completed geometry-calibration results so
that the link between the hardware configuration and the measured absolute
scales can be documented fully.

% ═══════════════════════════════════════════════════════════════════════════════
\begin{thebibliography}{99}

\bibitem{raffel2015bos}
M.~Raffel.
\newblock Background-oriented schlieren (BOS) techniques.
\newblock \emph{Experiments in Fluids}, 56:60, 2015.

\bibitem{chen2005cmm}
J.~Chen and J.~Katz.
\newblock Elimination of peak-locking error in PIV analysis using the
  correlation mapping method.
\newblock \emph{Measurement Science and Technology}, 16(8):1605--1618, 2005.

\bibitem{westerweel2005uno}
J.~Westerweel and F.~Scarano.
\newblock Universal outlier detection for PIV data.
\newblock \emph{Experiments in Fluids}, 39(6):1096--1100, 2005.

\bibitem{frankot1988}
R.~T.~Frankot and R.~Chellappa.
\newblock A method for enforcing integrability in shape from shading algorithms.
\newblock \emph{IEEE Transactions on Pattern Analysis and Machine Intelligence},
  10(4):439--451, 1988.

\bibitem{agrawal2006}
A.~Agrawal, R.~Chellappa, and R.~Raskar.
\newblock What is the range of surface reconstructions from a gradient field?
\newblock In \emph{Proceedings of ECCV}, pages 578--591, 2006.

\bibitem{liu2022uavbos}
X.~Liu, T.~Guo, P.~Zhang, Z.~Jia, and X.~Tong.
\newblock Extraction of a weak flow field for a multi-rotor unmanned aerial
  vehicle using high-speed background-oriented schlieren technology.
\newblock \emph{Sensors}, 22(1):43, 2022.

\end{thebibliography}

\end{document}
